% ============================================================
%  EXAMEN TEMA 4: FÍSICA — MEDICIONES, CINEMÁTICA Y GRÁFICAS
%  Curso Introductorio — 3er Año
%  Duración: 90 minutos
%  Temas: Magnitudes y unidades, conversión de unidades,
%         MRU y gráficas del movimiento
%  Autor: Ing. Gabriel Astudillo
%  Nota: enunciado limpio (sin pistas ni desarrollo). Las
%  soluciones están en la "Clave de Respuestas" al final.
% ============================================================

\documentclass[12pt, letterpaper]{article}

% ── Paquetes ─────────────────────────────────────────────────

\usepackage{mathwizards-guia}
\mwguia{Examen — Tema 4: Física: Mediciones, Cinemática y Gráficas}{Ing. Gabriel Astudillo $\cdot$ Curso 3er Año}

\begin{document}
% ============================================================

% ── Portada / Título ─────────────────────────────────────────
\mwportada{Examen del Tema 4}{Física: Mediciones, Cinemática y Gráficas}{Magnitudes, Conversión de Unidades, MRU y Gráficas}

\vspace{4pt}
\begin{cuadroinfo}
  \small
  \textbf{Nombre:} \rule{0.55\linewidth}{0.4pt} \quad \textbf{Fecha:} \rule{0.15\linewidth}{0.4pt} \\[6pt]
  \textbf{Tiempo disponible:} 90 minutos \quad$\bullet$\quad \textbf{Puntaje total:} 100 puntos \\[4pt]
  \textbf{Instrucciones:}
  \begin{enumerate}[leftmargin=*]
    \item Lee cada ejercicio con calma antes de resolverlo.
    \item Muestra \emph{todos} tus pasos, incluidos los factores de conversión.
    \item \textbf{No olvides las unidades} en tus resultados.
    \item Traza las gráficas en el plano cartesiano con ejes y escalas rotuladas.
    \item No se permite el uso de calculadora.
    \item Escribe con letra clara y revisa tu examen antes de entregarlo.
  \end{enumerate}
\end{cuadroinfo}

\vspace{8pt}

% ============================================================
\section{Fundamentos de la Medición (20 puntos)}
% ============================================================

\begin{enumerate}[leftmargin=*]
  \item \textbf{(6 pts)} Clasifica cada magnitud como escalar (E) o vectorial (V):
        \[
          a)\ \text{masa} \quad b)\ \text{fuerza} \quad c)\ \text{temperatura} \quad d)\ \text{desplazamiento} \quad e)\ \text{tiempo} \quad f)\ \text{velocidad}
        \]

  \item \textbf{(6 pts)} Completa la unidad del SI y su símbolo para cada magnitud:
        \[
          \text{longitud}, \quad \text{masa}, \quad \text{tiempo}, \quad \text{temperatura}, \quad \text{corriente eléctrica}, \quad \text{cantidad de sustancia}
        \]

  \item \textbf{(8 pts)}
        \begin{enumerate}[label=\alph*)]
          \item Explica la diferencia entre \textbf{rapidez} y \textbf{velocidad}, con un ejemplo cotidiano.
          \item Calcula el peso de un objeto de $25\;\text{kg}$ si $g = 9{,}8\;\text{m/s}^2$. \textquestiondown El peso es escalar o vectorial? Justifica.
        \end{enumerate}
\end{enumerate}

\newpage

% ============================================================
\section{Conversión de Unidades (25 puntos)}
% ============================================================

\begin{enumerate}[leftmargin=*]
  \item \textbf{(10 pts)} Convierte usando factores de conversión:
        \[
          a)\ 4{,}5\;\text{km} \to \text{m} \quad b)\ 3\;\text{h} \to \text{s} \quad c)\ 108\;\text{km/h} \to \text{m/s} \quad d)\ 15\;\text{m/s} \to \text{km/h} \quad e)\ 7{,}5\;\text{g/cm}^3 \to \text{kg/m}^3
        \]

  \item \textbf{(8 pts)} Convierte las siguientes temperaturas:
        \[
          a)\ 25\;^\circ\text{C} \to \text{K} \qquad b)\ 300\;\text{K} \to {}^\circ\text{C} \qquad c)\ 50\;^\circ\text{C} \to {}^\circ\text{F} \qquad d)\ 68\;^\circ\text{F} \to {}^\circ\text{C}
        \]

  \item \textbf{(7 pts)}
        \begin{enumerate}[label=\alph*)]
          \item La densidad del aluminio es $2\,700\;\text{kg/m}^3$. Exprésala en $\text{g/cm}^3$.
          \item La velocidad de la luz es $3 \times 10^8\;\text{m/s}$. Exprésala en $\text{km/h}$ usando notación científica.
        \end{enumerate}
\end{enumerate}

\newpage

% ============================================================
\section{Cinemática: Movimiento Rectilíneo Uniforme (30 puntos)}
% ============================================================

\begin{enumerate}[leftmargin=*]
  \item \textbf{(8 pts)} Resuelve:
        \begin{enumerate}[label=\alph*)]
          \item Un ciclista viaja a $12\;\text{m/s}$ durante $25\;\text{s}$. \textquestiondown Qué distancia recorre?
          \item Un auto recorre $180\;\text{km}$ en $2{,}5\;\text{h}$. \textquestiondown Cuál es su velocidad en km/h?
          \item \textquestiondown Cuánto tarda un corredor que va a $8\;\text{m/s}$ en recorrer $500\;\text{m}$?
        \end{enumerate}

  \item \textbf{(8 pts)} Un auto parte de la posición $x_0 = 80\;\text{m}$ y viaja a $-12\;\text{m/s}$ (sentido negativo). \textquestiondown Cuál es su posición después de $15\;\text{s}$? Interpreta el signo del resultado.

  \item \textbf{(7 pts)} Dos ciudades están separadas por $450\;\text{km}$. Dos autos parten \emph{simultáneamente} uno hacia el otro: el primero a $90\;\text{km/h}$ y el segundo a $60\;\text{km/h}$. \textquestiondown Después de cuántas horas se encuentran? \textquestiondown A qué distancia de la primera ciudad?

  \item \textbf{(7 pts)} Un auto $A$ sale de un punto a $50\;\text{km/h}$. Una hora y media después, el auto $B$ sale del mismo punto a $80\;\text{km/h}$ en la misma dirección. \textquestiondown Después de cuántas horas (desde que sale $B$) lo alcanza?
\end{enumerate}

\newpage

% ============================================================
\section{Gráficas del Movimiento (25 puntos)}
% ============================================================

\begin{enumerate}[leftmargin=*]
  \item \textbf{(8 pts)} En una gráfica $x$-$t$, un objeto se mueve en línea recta y pasa por los puntos $(0\;\text{s},\ 10\;\text{m})$ y $(5\;\text{s},\ 60\;\text{m})$. Calcula la velocidad del objeto y escribe la ecuación del movimiento.

  \item \textbf{(8 pts)} Un auto viaja a $24\;\text{m/s}$ durante $5\;\text{s}$, luego se detiene durante $4\;\text{s}$ y finalmente viaja a $-16\;\text{m/s}$ durante $3\;\text{s}$.
        \begin{enumerate}[label=\alph*)]
          \item Dibuja la gráfica $v$-$t$.
          \item Calcula el desplazamiento total usando el área bajo la curva.
        \end{enumerate}

  \item \textbf{(9 pts)} Dos móviles parten del mismo punto. El móvil $A$ viaja a $12\;\text{m/s}$. El móvil $B$ parte $3\;\text{s}$ después a $18\;\text{m/s}$ en la misma dirección. Determina, usando las ecuaciones, en qué instante y a qué distancia del origen $B$ alcanza a $A$.
\end{enumerate}

\newpage

% ============================================================
\ifmwclaves
  \section{Clave de Respuestas}
  % ============================================================

  \subsection*{Sección 1: Fundamentos de la Medición}

  \begin{enumerate}[leftmargin=*, label=\textbf{\arabic*.}]
    \item a) E \quad b) V \quad c) E \quad d) V \quad e) E \quad f) V

    \item Longitud: metro (m). Masa: kilogramo (kg). Tiempo: segundo (s). Temperatura: kelvin (K). Corriente eléctrica: amperio (A). Cantidad de sustancia: mol (mol).

    \item a) La \textbf{rapidez} es escalar: solo indica cuánto recorre por unidad de tiempo (siempre positiva). La \textbf{velocidad} es vectorial: además del módulo, tiene dirección y sentido (puede ser negativa). Ejemplo: un auto a $60\;\text{km/h}$ tiene esa rapidez, pero su velocidad puede ser $60\;\text{km/h}$ hacia el norte o hacia el sur.
          \\[2pt]
          b) Peso $= m\,g = 25 \times 9{,}8 = 245\;\text{N}$. El peso es \textbf{vectorial}: tiene módulo ($245\;\text{N}$), dirección (vertical) y sentido (hacia abajo, hacia el centro de la Tierra).
  \end{enumerate}

  \newpage

  \subsection*{Sección 2: Conversión de Unidades}

  \begin{enumerate}[leftmargin=*, label=\textbf{\arabic*.}]
    \item
          a) $4{,}5\;\text{km} \times \dfrac{1\,000\;\text{m}}{1\;\text{km}} = 4\,500\;\text{m}$.
          \quad b) $3\;\text{h} \times \dfrac{3\,600\;\text{s}}{1\;\text{h}} = 10\,800\;\text{s}$.
          \quad c) $108 \div 3{,}6 = 30\;\text{m/s}$.
          \quad d) $15 \times 3{,}6 = 54\;\text{km/h}$.
          \quad e) $7{,}5\;\text{g/cm}^3 \times 1\,000 = 7\,500\;\text{kg/m}^3$.

    \item a) $T_K = 25 + 273{,}15 = 298{,}15\;\text{K}$.
          \quad b) $T_C = 300 - 273{,}15 = 26{,}85\;^\circ\text{C}$.
          \quad c) $T_F = \frac{9}{5}(50) + 32 = 90 + 32 = 122\;^\circ\text{F}$.
          \quad d) $T_C = \frac{5}{9}(68 - 32) = \frac{5}{9}(36) = 20\;^\circ\text{C}$.

    \item a) $1\;\text{kg/m}^3 = 0{,}001\;\text{g/cm}^3$, entonces $2\,700\;\text{kg/m}^3 = 2{,}7\;\text{g/cm}^3$.
          \quad b) $3 \times 10^8\;\text{m/s} \times 3{,}6 = 10{,}8 \times 10^8 = 1{,}08 \times 10^9\;\text{km/h}$.
  \end{enumerate}

  \newpage

  \subsection*{Sección 3: Cinemática: Movimiento Rectilíneo Uniforme}

  \begin{enumerate}[leftmargin=*, label=\textbf{\arabic*.}]
    \item a) $d = v\,t = 12 \times 25 = 300\;\text{m}$. \quad
          b) $v = \dfrac{180}{2{,}5} = 72\;\text{km/h}$. \quad
          c) $t = \dfrac{500}{8} = 62{,}5\;\text{s}$.

    \item $x_f = x_0 + v\,t = 80 + (-12)(15) = 80 - 180 = -100\;\text{m}$.
          El signo negativo indica que el auto pasó por el origen y ahora se encuentra $100\;\text{m}$ al otro lado (en el sentido negativo).

    \item Velocidad de acercamiento: $90 + 60 = 150\;\text{km/h}$. \quad $t = \dfrac{450}{150} = 3\;\text{h}$.
          Distancia desde la primera ciudad: $90 \times 3 = 270\;\text{km}$.

    \item Cuando $B$ sale, $A$ lleva una ventaja de $50 \times 1{,}5 = 75\;\text{km}$.
          Velocidad de acercamiento: $80 - 50 = 30\;\text{km/h}$. \quad $t = \dfrac{75}{30} = 2{,}5\;\text{h}$ (desde que sale $B$).
  \end{enumerate}

  \newpage

  \subsection*{Sección 4: Gráficas del Movimiento}

  \begin{enumerate}[leftmargin=*, label=\textbf{\arabic*.}]
    \item $v = \dfrac{60 - 10}{5 - 0} = \dfrac{50}{5} = 10\;\text{m/s}$.
          Con $x_0 = 10\;\text{m}$: ecuación $x = 10 + 10t$ (con $x$ en metros y $t$ en segundos).

    \item a) Gráfica $v$-$t$: rectángulo positivo en $v = 24\;\text{m/s}$ de $t = 0$ a $5\;\text{s}$; luego $v = 0$ de $5$ a $9\;\text{s}$; y un rectángulo negativo en $v = -16\;\text{m/s}$ de $9$ a $12\;\text{s}$.
          \\[2pt]
          b) $\Delta x = (24 \times 5) + (0 \times 4) + (-16 \times 3) = 120 + 0 - 48 = 72\;\text{m}$.

    \item Ecuaciones: $x_A = 12t$ y $x_B = 18(t - 3)$ para $t \geq 3$.
          Igualando: $12t = 18(t - 3) \Rightarrow 12t = 18t - 54 \Rightarrow 6t = 54 \Rightarrow t = 9\;\text{s}$.
          Distancia: $x = 12(9) = 108\;\text{m}$ del origen.
  \end{enumerate}

\fi
\end{document}
